
\documentclass{article}
\usepackage{colortbl}
\usepackage{makecell}
\usepackage{multirow}
\usepackage{supertabular}

\begin{document}

\newcounter{utterance}

\centering \large Interaction Transcript for game `dond', experiment `semi\_en', episode 1 with Qwen3{-}30B{-}A3B{-}t0.0.
\vspace{24pt}

{ \footnotesize  \setcounter{utterance}{1}
\setlength{\tabcolsep}{0pt}
\begin{supertabular}{c@{$\;$}|p{.15\linewidth}@{}p{.15\linewidth}p{.15\linewidth}p{.15\linewidth}p{.15\linewidth}p{.15\linewidth}}
    \# & $\;$A & \multicolumn{4}{c}{Game Master} & $\;\:$B\\
    \hline

    \theutterance \stepcounter{utterance}  
    & & \multicolumn{4}{p{0.6\linewidth}}{
        \cellcolor[rgb]{0.9,0.9,0.9}{
            \makecell[{{p{\linewidth}}}]{
                \texttt{\tiny{[P1$\langle$GM]}}
                \texttt{You are playing a negotiation game in which you have to agree on how to divide a set of items among you and another player.} \\
\\ 
\texttt{Rules:} \\
\texttt{(a) You and the other player are given a set of items. Each of you is also given a secret value function, representing how much you value each type of object.} \\
\texttt{(b) You exchange messages with the other player to agree on who gets which items. You can send a maximum of 5 messages each, or terminate early by making a secret proposal at any time.} \\
\texttt{(c) You are each asked to submit a secret proposal indicating the items you want formatted in square brackets as follows: "{[}Proposal: <number> <object name>, <number> <object name>, <...>{]}"} \\
\texttt{(d) If your proposals are complementary, i.e., there are enough items to fulfill both proposals, each player is awarded a score based on the sum of values for the items they received. Otherwise, both of you get zero points.} \\
\texttt{(e) Your goal is to maximize the score you receive.} \\
\\ 
\texttt{Let us start.} \\
\\ 
\texttt{The set of available items is:} \\
\\ 
\texttt{4 magnets, 1 shovel, 1 flower} \\
\\ 
\texttt{Your secret value function is:} \\
\\ 
\texttt{magnet: 0, shovel: 2, flower: 8} \\
\\ 
\texttt{IMPORTANT: Your messages, unless it is the secret proposal, are directly transmitted to the other player, so do not include any response to the rules or text announcing your message. To make a secret proposal, use the indicated format. Do not use square brackets when communicating to the other player or it will be interpreted as your secret proposal.} \\
            }
        }
    }
    & \\ \\

    \theutterance \stepcounter{utterance}  
    & \multicolumn{4}{p{0.6\linewidth}}{
        \cellcolor[rgb]{0.9,0.9,0.9}{
            \makecell[{{p{\linewidth}}}]{
                \texttt{\tiny{[P1$\rangle$GM]}}
                \texttt{<think>} \\
\texttt{Okay, let's see. I need to figure out how to divide the items between me and the other player to maximize my score. The items are 4 magnets, 1 shovel, and 1 flower. My value function is magnets: 0, shovel: 2, flower: 8. So I don't value magnets at all, but the flower is really valuable to me. The other player's values are secret, so I don't know what they prefer.} \\
\\ 
\texttt{The rules say that if our proposals are complementary, meaning there are enough items to fulfill both, then we get points based on our own values. Otherwise, both get zero. So I need to make a proposal that the other player can accept, and also ensure that the items can be divided without conflict.} \\
\\ 
\texttt{Since I don't know the other player's values, I have to make an educated guess. But maybe I can think about what they might want. The flower is worth 8 to me, so maybe they value it less or more? If they value the flower highly, they might want it. But if they don't, maybe they'd prefer the shovel or magnets. But magnets are worth 0 to me, so maybe they have some value for them?} \\
\\ 
\texttt{Wait, but the other player's values are secret. So I can't be sure. But the key is to make a proposal that the other player would accept, and that doesn't conflict with their possible needs.} \\
\\ 
\texttt{Let me think: If I propose to take the flower and the shovel, but the other player might want the flower. But if I take the flower, they might not get it. Alternatively, if I propose to take the flower, and the other player takes the shovel and magnets, that could work. But how to communicate that?} \\
\\ 
\texttt{But since I can send messages, maybe I can try to negotiate. But I can only send 5 messages. Let me think about possible strategies.} \\
\\ 
\texttt{First, my own proposal: I want the flower (value 8) and the shovel (value 2). But if I take both, that's 10 points. But if the other player wants the flower, they might not accept that. Alternatively, if I take the flower and the other player takes the shovel and magnets, that's possible. But how to get them to agree?} \\
\\ 
\texttt{Alternatively, maybe I can propose to take the flower and the other player takes the shovel and magnets. But how to make sure that the} \\
            }
        }
    }
    & & \\ \\

    \theutterance \stepcounter{utterance}  
    & & & \multicolumn{2}{p{0.3\linewidth}}{
        \cellcolor[rgb]{0.9,0.9,0.9}{
            \makecell[{{p{\linewidth}}}]{
                \texttt{\tiny{[GM$|$GM]}}
                \texttt{continue} \\
            }
        }
    }
    & & \\ \\

    \theutterance \stepcounter{utterance}  
    & & & \multicolumn{2}{p{0.3\linewidth}}{
        \cellcolor[rgb]{0.9,0.9,0.9}{
            \makecell[{{p{\linewidth}}}]{
                \texttt{\tiny{[GM$|$GM]}}
                \texttt{<think>} \\
\texttt{Okay, let's see. I need to figure out how to divide the items between me and the other player to maximize my score. The items are 4 magnets, 1 shovel, and 1 flower. My value function is magnets: 0, shovel: 2, flower: 8. So I don't value magnets at all, but the flower is really valuable to me. The other player's values are secret, so I don't know what they prefer.} \\
\\ 
\texttt{The rules say that if our proposals are complementary, meaning there are enough items to fulfill both, then we get points based on our own values. Otherwise, both get zero. So I need to make a proposal that the other player can accept, and also ensure that the items can be divided without conflict.} \\
\\ 
\texttt{Since I don't know the other player's values, I have to make an educated guess. But maybe I can think about what they might want. The flower is worth 8 to me, so maybe they value it less or more? If they value the flower highly, they might want it. But if they don't, maybe they'd prefer the shovel or magnets. But magnets are worth 0 to me, so maybe they have some value for them?} \\
\\ 
\texttt{Wait, but the other player's values are secret. So I can't be sure. But the key is to make a proposal that the other player would accept, and that doesn't conflict with their possible needs.} \\
\\ 
\texttt{Let me think: If I propose to take the flower and the shovel, but the other player might want the flower. But if I take the flower, they might not get it. Alternatively, if I propose to take the flower, and the other player takes the shovel and magnets, that could work. But how to communicate that?} \\
\\ 
\texttt{But since I can send messages, maybe I can try to negotiate. But I can only send 5 messages. Let me think about possible strategies.} \\
\\ 
\texttt{First, my own proposal: I want the flower (value 8) and the shovel (value 2). But if I take both, that's 10 points. But if the other player wants the flower, they might not accept that. Alternatively, if I take the flower and the other player takes the shovel and magnets, that's possible. But how to get them to agree?} \\
\\ 
\texttt{Alternatively, maybe I can propose to take the flower and the other player takes the shovel and magnets. But how to make sure that the} \\
            }
        }
    }
    & & \\ \\

    \theutterance \stepcounter{utterance}  
    & & \multicolumn{4}{p{0.6\linewidth}}{
        \cellcolor[rgb]{0.9,0.9,0.9}{
            \makecell[{{p{\linewidth}}}]{
                \texttt{\tiny{[GM$\rangle$P2]}}
                \texttt{You are playing a negotiation game in which you have to agree on how to divide a set of items among you and another player.} \\
\\ 
\texttt{Rules:} \\
\texttt{(a) You and the other player are given a set of items. Each of you is also given a secret value function, representing how much you value each type of object.} \\
\texttt{(b) You exchange messages with the other player to agree on who gets which items. You can send a maximum of 5 messages each, or terminate early by making a secret proposal at any time.} \\
\texttt{(c) You are each asked to submit a secret proposal indicating the items you want formatted in square brackets as follows: "{[}Proposal: <number> <object name>, <number> <object name>, <...>{]}"} \\
\texttt{(d) If your proposals are complementary, i.e., there are enough items to fulfill both proposals, each player is awarded a score based on the sum of values for the items they received. Otherwise, both of you get zero points.} \\
\texttt{(e) Your goal is to maximize the score you receive.} \\
\\ 
\texttt{Let us start.} \\
\\ 
\texttt{The set of available items is:} \\
\\ 
\texttt{4 magnets, 1 shovel, 1 flower} \\
\\ 
\texttt{Your secret value function is:} \\
\\ 
\texttt{magnet: 1, shovel: 4, flower: 2} \\
\\ 
\texttt{IMPORTANT: Your messages, unless it is the secret proposal, are directly transmitted to the other player, so do not include any response to the rules or text announcing your message. To make a secret proposal, use the indicated format. Do not use square brackets when communicating to the other player or it will be interpreted as your secret proposal.} \\
\\ 
\\ 
\texttt{<think>} \\
\texttt{Okay, let's see. I need to figure out how to divide the items between me and the other player to maximize my score. The items are 4 magnets, 1 shovel, and 1 flower. My value function is magnets: 0, shovel: 2, flower: 8. So I don't value magnets at all, but the flower is really valuable to me. The other player's values are secret, so I don't know what they prefer.} \\
\\ 
\texttt{The rules say that if our proposals are complementary, meaning there are enough items to fulfill both, then we get points based on our own values. Otherwise, both get zero. So I need to make a proposal that the other player can accept, and also ensure that the items can be divided without conflict.} \\
\\ 
\texttt{Since I don't know the other player's values, I have to make an educated guess. But maybe I can think about what they might want. The flower is worth 8 to me, so maybe they value it less or more? If they value the flower highly, they might want it. But if they don't, maybe they'd prefer the shovel or magnets. But magnets are worth 0 to me, so maybe they have some value for them?} \\
\\ 
\texttt{Wait, but the other player's values are secret. So I can't be sure. But the key is to make a proposal that the other player would accept, and that doesn't conflict with their possible needs.} \\
\\ 
\texttt{Let me think: If I propose to take the flower and the shovel, but the other player might want the flower. But if I take the flower, they might not get it. Alternatively, if I propose to take the flower, and the other player takes the shovel and magnets, that could work. But how to communicate that?} \\
\\ 
\texttt{But since I can send messages, maybe I can try to negotiate. But I can only send 5 messages. Let me think about possible strategies.} \\
\\ 
\texttt{First, my own proposal: I want the flower (value 8) and the shovel (value 2). But if I take both, that's 10 points. But if the other player wants the flower, they might not accept that. Alternatively, if I take the flower and the other player takes the shovel and magnets, that's possible. But how to get them to agree?} \\
\\ 
\texttt{Alternatively, maybe I can propose to take the flower and the other player takes the shovel and magnets. But how to make sure that the} \\
            }
        }
    }
    & \\ \\

    \theutterance \stepcounter{utterance}  
    & & & \multicolumn{4}{p{0.6\linewidth}}{
        \cellcolor[rgb]{0.9,0.9,0.9}{
            \makecell[{{p{\linewidth}}}]{
                \texttt{\tiny{[GM$\langle$P2]}}
                \texttt{<think>} \\
\texttt{Okay, let's try to figure out the best approach here. The goal is to maximize my score, which depends on the items I get based on my value function. My values are: magnets (0), shovel (2), flower (8). So I really want the flower, then the shovel, and magnets are worthless to me. The other player's values are unknown, but I need to make a proposal that they can accept without conflict.} \\
\\ 
\texttt{First, I need to consider that if my proposal and theirs are complementary, we both get points. Otherwise, zero. So I need to make a proposal that the other player can agree with. But since I don't know their values, I have to make an assumption.} \\
\\ 
\texttt{If I take the flower, which is high value for me, but maybe the other player also wants it. If I take the flower, they might not get it, but if they have a high value for it, they might not accept my proposal. Alternatively, if I don't take the flower, maybe they can take it, and I can take the shovel and magnets. But since magnets are worth 0 to me, I don't care about them. So maybe I should focus on getting the flower and the shovel.} \\
\\ 
\texttt{But how to ensure that the other player's proposal doesn't conflict. Let's think about possible scenarios.} \\
\\ 
\texttt{If I propose to take the flower and the shovel, then the other player would need to take the magnets. But if they value the magnets, that's okay. But if they don't, maybe they would want something else. However, since there's only one shovel and one flower, if I take both, the other player can only take the magnets. But since there are 4 magnets, they can take all of them. But since magnets are worth 0 to me, that's fine. But the other player might not want magnets. But I don't know.} \\
\\ 
\texttt{Alternatively, maybe I should propose to take the flower and the other player takes the shovel and magnets. But how to communicate that. Since I can send messages, maybe I can suggest that. But the problem is that the other player's values are unknown. So I need to make a proposal that is likely to be compatible.} \\
\\ 
\texttt{Another angle: The total items are 4 magnets, 1 shovel, 1 flower. If I take the flower, that's 8 points. If I take the shovel, that's 2 points. If I} \\
            }
        }
    }
     \\ \\

    \theutterance \stepcounter{utterance}  
    & & & \multicolumn{2}{p{0.3\linewidth}}{
        \cellcolor[rgb]{0.9,0.9,0.9}{
            \makecell[{{p{\linewidth}}}]{
                \texttt{\tiny{[GM$|$GM]}}
                \texttt{continue} \\
            }
        }
    }
    & & \\ \\

    \theutterance \stepcounter{utterance}  
    & & & \multicolumn{2}{p{0.3\linewidth}}{
        \cellcolor[rgb]{0.9,0.9,0.9}{
            \makecell[{{p{\linewidth}}}]{
                \texttt{\tiny{[GM$|$GM]}}
                \texttt{<think>} \\
\texttt{Okay, let's try to figure out the best approach here. The goal is to maximize my score, which depends on the items I get based on my value function. My values are: magnets (0), shovel (2), flower (8). So I really want the flower, then the shovel, and magnets are worthless to me. The other player's values are unknown, but I need to make a proposal that they can accept without conflict.} \\
\\ 
\texttt{First, I need to consider that if my proposal and theirs are complementary, we both get points. Otherwise, zero. So I need to make a proposal that the other player can agree with. But since I don't know their values, I have to make an assumption.} \\
\\ 
\texttt{If I take the flower, which is high value for me, but maybe the other player also wants it. If I take the flower, they might not get it, but if they have a high value for it, they might not accept my proposal. Alternatively, if I don't take the flower, maybe they can take it, and I can take the shovel and magnets. But since magnets are worth 0 to me, I don't care about them. So maybe I should focus on getting the flower and the shovel.} \\
\\ 
\texttt{But how to ensure that the other player's proposal doesn't conflict. Let's think about possible scenarios.} \\
\\ 
\texttt{If I propose to take the flower and the shovel, then the other player would need to take the magnets. But if they value the magnets, that's okay. But if they don't, maybe they would want something else. However, since there's only one shovel and one flower, if I take both, the other player can only take the magnets. But since there are 4 magnets, they can take all of them. But since magnets are worth 0 to me, that's fine. But the other player might not want magnets. But I don't know.} \\
\\ 
\texttt{Alternatively, maybe I should propose to take the flower and the other player takes the shovel and magnets. But how to communicate that. Since I can send messages, maybe I can suggest that. But the problem is that the other player's values are unknown. So I need to make a proposal that is likely to be compatible.} \\
\\ 
\texttt{Another angle: The total items are 4 magnets, 1 shovel, 1 flower. If I take the flower, that's 8 points. If I take the shovel, that's 2 points. If I} \\
            }
        }
    }
    & & \\ \\

    \theutterance \stepcounter{utterance}  
    & & \multicolumn{4}{p{0.6\linewidth}}{
        \cellcolor[rgb]{0.9,0.9,0.9}{
            \makecell[{{p{\linewidth}}}]{
                \texttt{\tiny{[P1$\langle$GM]}}
                \texttt{<think>} \\
\texttt{Okay, let's try to figure out the best approach here. The goal is to maximize my score, which depends on the items I get based on my value function. My values are: magnets (0), shovel (2), flower (8). So I really want the flower, then the shovel, and magnets are worthless to me. The other player's values are unknown, but I need to make a proposal that they can accept without conflict.} \\
\\ 
\texttt{First, I need to consider that if my proposal and theirs are complementary, we both get points. Otherwise, zero. So I need to make a proposal that the other player can agree with. But since I don't know their values, I have to make an assumption.} \\
\\ 
\texttt{If I take the flower, which is high value for me, but maybe the other player also wants it. If I take the flower, they might not get it, but if they have a high value for it, they might not accept my proposal. Alternatively, if I don't take the flower, maybe they can take it, and I can take the shovel and magnets. But since magnets are worth 0 to me, I don't care about them. So maybe I should focus on getting the flower and the shovel.} \\
\\ 
\texttt{But how to ensure that the other player's proposal doesn't conflict. Let's think about possible scenarios.} \\
\\ 
\texttt{If I propose to take the flower and the shovel, then the other player would need to take the magnets. But if they value the magnets, that's okay. But if they don't, maybe they would want something else. However, since there's only one shovel and one flower, if I take both, the other player can only take the magnets. But since there are 4 magnets, they can take all of them. But since magnets are worth 0 to me, that's fine. But the other player might not want magnets. But I don't know.} \\
\\ 
\texttt{Alternatively, maybe I should propose to take the flower and the other player takes the shovel and magnets. But how to communicate that. Since I can send messages, maybe I can suggest that. But the problem is that the other player's values are unknown. So I need to make a proposal that is likely to be compatible.} \\
\\ 
\texttt{Another angle: The total items are 4 magnets, 1 shovel, 1 flower. If I take the flower, that's 8 points. If I take the shovel, that's 2 points. If I} \\
            }
        }
    }
    & \\ \\

    \theutterance \stepcounter{utterance}  
    & \multicolumn{4}{p{0.6\linewidth}}{
        \cellcolor[rgb]{0.9,0.9,0.9}{
            \makecell[{{p{\linewidth}}}]{
                \texttt{\tiny{[P1$\rangle$GM]}}
                \texttt{<think>} \\
\texttt{Okay, let's try to structure this. My value function is clear: flower is 8, shovel 2, magnets 0. The other player's values are unknown, but I need to make a proposal that they can accept. Since I can send messages, maybe I can try to negotiate.} \\
\\ 
\texttt{First, I should consider what the other player might want. If they value the flower, they might want it. If they don't, maybe they prefer the shovel or magnets. But since I don't know, I have to make a proposal that's flexible.} \\
\\ 
\texttt{If I propose to take the flower, that's 8 points for me. But if the other player also wants the flower, our proposals would conflict. So maybe I should propose to take the flower and the shovel, but that might not be compatible if they want either. Alternatively, maybe I can propose to take the flower and let them take the rest. But how to ensure that.} \\
\\ 
\texttt{Wait, the key is that the proposals need to be complementary. So if I take the flower, and they take the shovel and magnets, that works. But how to get them to agree. Since I can send messages, maybe I can suggest that. But the messages are directly transmitted, so I need to be careful.} \\
\\ 
\texttt{Alternatively, maybe I should start by sending a message asking what they want. But since I can only send 5 messages, maybe I should be strategic.} \\
\\ 
\texttt{But the problem is that I can't ask them directly; I have to make a proposal. So perhaps the best approach is to propose a division that gives me the flower and the shovel, and the other player gets the magnets. But since there are 4 magnets, that's a lot, but they might not care. Alternatively, maybe the other player would prefer the shovel and magnets, and I take the flower.} \\
\\ 
\texttt{But how to ensure that. Let me think: If I propose to take the flower, and the other player proposes to take the shovel and magnets, that's complementary. So my proposal would be {[}Proposal: 1 flower{]}, and theirs would be {[}Proposal: 1 shovel, 4 magnets{]}. That works. But how to get them to agree to that.} \\
\\ 
\texttt{But since I can't communicate their proposal, I have to make my own proposal in a way that they can accept. So maybe I should propose to take the flower, and in a message, suggest that they take the shovel and magnets. But the messages are} \\
            }
        }
    }
    & & \\ \\

    \theutterance \stepcounter{utterance}  
    & & & \multicolumn{2}{p{0.3\linewidth}}{
        \cellcolor[rgb]{0.9,0.9,0.9}{
            \makecell[{{p{\linewidth}}}]{
                \texttt{\tiny{[GM$|$GM]}}
                \texttt{proposal} \\
            }
        }
    }
    & & \\ \\

    \theutterance \stepcounter{utterance}  
    & & & \multicolumn{2}{p{0.3\linewidth}}{
        \cellcolor[rgb]{0.9,0.9,0.9}{
            \makecell[{{p{\linewidth}}}]{
                \texttt{\tiny{[GM$|$GM]}}
                [0, 0, 1]
            }
        }
    }
    & & \\ \\

    \theutterance \stepcounter{utterance}  
    & & \multicolumn{4}{p{0.6\linewidth}}{
        \cellcolor[rgb]{0.9,0.9,0.9}{
            \makecell[{{p{\linewidth}}}]{
                \texttt{\tiny{[GM$\rangle$P2]}}
                \texttt{The other player has submitted a secret proposal. It is now your turn to submit a proposal of your own.} \\
            }
        }
    }
    & \\ \\

    \theutterance \stepcounter{utterance}  
    & & & \multicolumn{4}{p{0.6\linewidth}}{
        \cellcolor[rgb]{0.9,0.9,0.9}{
            \makecell[{{p{\linewidth}}}]{
                \texttt{\tiny{[GM$\langle$P2]}}
                \texttt{<think>} \\
\texttt{Okay, the other player has already submitted their secret proposal. Now it's my turn to submit mine. I need to figure out what their proposal might be and adjust mine accordingly.} \\
\\ 
\texttt{Since I don't know their values, I have to make an educated guess. Let's think about possible scenarios. The other player might value the flower, shovel, or magnets differently. But since I know my own values, I can try to find a proposal that complements theirs.} \\
\\ 
\texttt{If I assume that the other player wants the flower, then I should avoid conflicting with that. But if they don't, maybe they want the shovel or magnets. However, since I don't have information, I need to make a proposal that maximizes my chances.} \\
\\ 
\texttt{My value function is magnets: 0, shovel: 2, flower: 8. So I want the flower and shovel. If I can get both, that's 10 points. But if the other player also wants the flower, they might not accept my proposal. Alternatively, if I take the flower and they take the shovel and magnets, that could work.} \\
\\ 
\texttt{But since I don't know their proposal, I need to make a proposal that is likely to be complementary. Let's consider that the other player might have a proposal that includes the shovel or magnets. If I take the flower and the shovel, that would take both the flower and shovel, leaving the magnets for them. But if they want the shovel, that would conflict. Alternatively, if I take the flower and they take the shovel and magnets, that's complementary.} \\
\\ 
\texttt{But how to ensure that. Maybe I should propose to take the flower and the other player takes the shovel and magnets. But how to communicate that. Since I can only send messages, but the other player already submitted their proposal. Wait, the user said "The other player has submitted a secret proposal. It is now your turn to submit a proposal of your own." So I need to submit my proposal now, based on their proposal.} \\
\\ 
\texttt{But since I don't know their proposal, I have to make an assumption. Maybe the other player's proposal is something like taking the shovel and magnets. If that's the case, then my proposal to take the flower would be complementary. Alternatively, if their proposal is to take the flower, then my proposal to take the flower would conflict.} \\
\\ 
\texttt{But since I don't know, I have to make a choice. Given that the flower is worth 8 to me} \\
            }
        }
    }
     \\ \\

    \theutterance \stepcounter{utterance}  
    & & & \multicolumn{2}{p{0.3\linewidth}}{
        \cellcolor[rgb]{0.9,0.9,0.9}{
            \makecell[{{p{\linewidth}}}]{
                \texttt{\tiny{[GM$|$GM]}}
                \texttt{continue} \\
            }
        }
    }
    & & \\ \\

    \theutterance \stepcounter{utterance}  
    & & & \multicolumn{2}{p{0.3\linewidth}}{
        \cellcolor[rgb]{0.9,0.9,0.9}{
            \makecell[{{p{\linewidth}}}]{
                \texttt{\tiny{[GM$|$GM]}}
                \texttt{player was instructed to make a proposal} \\
            }
        }
    }
    & & \\ \\

\end{supertabular}
}

\end{document}
